% !TeX program = lualatex
% Compile twice: lualatex 122.tex
% All references and prose are embedded; no BibTeX database or external inputs.
\documentclass[11pt,a4paper]{article}
\usepackage[margin=24mm,headheight=14pt]{geometry}
\usepackage{fontspec}
\setmainfont{Latin Modern Roman}
\setsansfont{Latin Modern Sans}
\setmonofont{Latin Modern Mono}
\usepackage{amsmath,amssymb,mathtools}
\usepackage{unicode-math}
\setmathfont{Latin Modern Math}
\usepackage{microtype}
\usepackage{enumitem,xurl,needspace}
\usepackage[dvipsnames]{xcolor}
\usepackage{hyperref}
\hypersetup{unicode=true,colorlinks=true,linkcolor=MidnightBlue,urlcolor=MidnightBlue,
 pdftitle={500 Mathematical Problem Statements},pdfauthor={Source-annotated compilation},bookmarksnumbered=true}
\usepackage{bookmark}
\usepackage{tocloft}
\setlength{\cftsecnumwidth}{3em}
\cftsetpnumwidth{2.5em}
\cftsetrmarg{4em}
\usepackage{titlesec}
\titleformat{\section}{\Large\bfseries\raggedright}{\thesection}{0.7em}{}
\usepackage{fancyhdr}
\pagestyle{fancy}\fancyhf{}
\fancyhead[L]{\small\sffamily Mathematical problem statements}
\fancyhead[R]{\small Source edition 22}
\fancyfoot[C]{\thepage}
\setlength{\parindent}{0pt}
\setlength{\parskip}{5pt plus 1pt}
\setlength{\emergencystretch}{3em}
\setcounter{tocdepth}{1}
\setcounter{secnumdepth}{1}
\setlist[itemize]{leftmargin=3em,itemsep=5pt,topsep=4pt,parsep=0pt}
\allowdisplaybreaks
\newcommand{\N}{\mathbb{N}}
\newcommand{\Z}{\mathbb{Z}}
\newcommand{\Q}{\mathbb{Q}}
\newcommand{\R}{\mathbb{R}}
\newcommand{\C}{\mathbb{C}}
\newcommand{\F}{\mathbb{F}}
\newcommand{\A}{\mathbb{A}}
\newcommand{\PP}{\mathbb P}
\newcommand{\E}{\mathbb{E}}
\newcommand{\eps}{\varepsilon}
\newcommand{\dd}{\,\mathrm d}
\newcommand{\sref}[2]{\hyperlink{source:#1:#2}{[S#2]}}
\newcommand{\eref}[2]{\hyperlink{extra:#1:#2}{[E#2]}}
% Turn the next switch off for a shorter reading copy. The quotations preserve
% every exactTarget field, including qualifications omitted from a short summary.
\newif\ifshowcatalogue
\showcataloguetrue
\newcommand{\cataloguescope}[1]{\ifshowcatalogue
 \paragraph{Exact scope in the supplied catalogue.}
 \begin{quote}\small #1\end{quote}\fi}
\hypersetup{pdftitle={Problem 122 - QMA versus QCMA},pdfauthor={Open Problems Math research notebook}}
\fancyhead[L]{\small\sffamily Open Problems Math \textbar\ Research notebook}
\fancyhead[R]{\small\sffamily 122 \textbar\ 27 September 2026}
\numberwithin{equation}{section}
\begin{document}
\setcounter{section}{121}
% ================================================================
% problemId: problem.qma-versus-qcma-problem
% source releaseRank: 122; source releaseStatus: open_with_solved_subcases
% exactTarget SHA-256: 80f2154356ed3603ca1ecf153dbf2005e8eed03d0746269c5c5dfeb7c27bb783
\clearpage
\section{\texorpdfstring{QMA versus QCMA Problem}{QMA versus QCMA Problem}}\label{122}
\subsection{Definitions and mathematical statement}
A QMA verifier is a uniform polynomial-time quantum computation receiving a polynomial-size quantum witness. A QCMA verifier has the same quantum computational power but receives only a polynomial-length classical bit string. Both use completeness at least $2/3$ and soundness at most $1/3$, with soundness quantified over every allowed witness on no-inputs. The question is
\[
 \mathsf{QMA}\stackrel{?}{=}\mathsf{QCMA}.
\]
The catalogue states a language formulation; a promise-problem version should be labeled separately rather than silently substituted. An oracle separation does not determine the ordinary unrelativized relation.
\par\smallskip\textit{Formulation and concept references:} \sref{122}{1}.
\cataloguescope{Determine whether QMA=QCMA: does every language decidable by a polynomial-time quantum verifier using a polynomial-size quantum witness also admit a polynomial-time quantum verification protocol using only a polynomial-size classical witness?}
\subsection{Short English statement}
Can quantum evidence always be replaced by classical evidence without sacrificing efficient quantum verification?
% BEGIN RESEARCH ADDITION 122
\subsection{Research attempt}

\paragraph{Current frontier.}
The cited Bostanci--Haferkamp--Nirkhe--Zhandry manuscript proves a
\emph{classical-oracle} separation, in a language formulation, and
explains why witness-state complexity alone does not settle the
unrelativized question.  Theorem 1.1 and the accompanying discussion
were checked in its January 2026 revision. \sref{122}{1}
The April 2026 revision of Bostanci--Huang--Vaikuntanathan gives another
classical-oracle separation using coding ideas. \eref{122}{1}
Neither statement determines the relation of ordinary QMA and QCMA.
We keep the supplied language convention and do not turn a promise-only
spectral problem into a total language without a valid construction.

\paragraph{Attack routes.}
To prove equality, one could compress an accepting quantum witness into
a classical description of a preparable state.  This needs existence
of a succinct accepting state, not succinctness of every quantum state.
To prove separation, one could try to force highly complex witnesses,
but must also exclude a different QCMA verifier for the same language.
A third possibility is to remove the oracle from the recent constructions;
it requires an efficiently computable explicit object retaining the
adversary lower bound, rather than an oracle chosen by diagonalization.
The compression route is developed here to quantify one sufficient
notion of classical describability and test its limitations.

\paragraph{Attempt I: acceptance operators and a compression lemma.}
Fix an input $x$ of length $n\geq2$.  A verifier acting on $m=m(n)$
witness qubits induces a positive operator $0\leq A_x\leq I$, with
acceptance probability $\langle\psi,A_x\psi\rangle$ for a pure witness.
Maximizing over density matrices gives $\|A_x\|$: diagonalize $A_x$,
and observe that the expectation is a convex combination of its
eigenvalues.  Thus completeness is an eigenvalue-existence assertion,
while soundness quantifies over the entire unit sphere.

Suppose that for every yes-input there is an accepting pure witness
$\psi$, and a polynomial-size quantum circuit $U$, such that the
probabilities
\[
 p_z=|\langle z,U^*\psi\rangle|^2
 \quad(z\in\{0,1\}^m)
\]
have Shannon entropy $\sum_zp_z\log_2(1/p_z)\leq C\log_2 n$, where
$C\geq1$ and the polynomial size bounds are fixed for the language
(enlarging $C$ if necessary).  The
circuit $U$ need not be computable from $x$; it may be part of the
classical witness.

\textbf{Entropy-compression lemma.} Under this extra hypothesis the
language has a QCMA verifier.

Take $\delta=1/576$ and let
$B=\{z:p_z\geq n^{-C/\delta}\}$.  There are at most
$n^{C/\delta}$ such strings.  For $z\notin B$,
$\log_2(1/p_z)>(C/\delta)\log_2 n$, so the entropy bound implies
\[
 \sum_{z\notin B}p_z\leq\delta.
\]
The normalized truncation $\phi$ of $U^*\psi$ to $B$ satisfies
$|\langle\phi,U^*\psi\rangle|^2\geq1-\delta$.  The trace distance
between these pure states is at most $\sqrt\delta=1/24$.  Therefore
\[
 \langle U\phi,A_xU\phi\rangle\geq\frac23-\frac1{24}.
\]
A classical description can list the support strings and their complex
amplitudes rounded to inverse-polynomial accuracy.  If there are $s$
entries, rounding each real and imaginary part within $O(1/s)$ and
renormalizing makes the Euclidean error a fixed small constant; choosing
more bits makes the resulting trace distance below $1/24$.
The bit count is $O(s(m+\log s))$, hence polynomial in $n$.

Such a sparse state can be prepared with polynomial overhead: prepare
a register indexing its $s$ entries by successive controlled rotations,
compute the listed support string reversibly, and erase the index using
the inverse lookup on the distinct strings.  Rotations and the supplied
circuit can be approximated over a fixed universal gate set within the
chosen total error.  The quantum verifier checks the syntax and size
bounds of the description, performs this preparation, and runs the
original verifier.  Its yes-acceptance probability is at least $7/12$.
On no-inputs \emph{every} resulting physical state is accepted with
probability at most $1/3$, whether or not the claimed amplitudes or basis
have any relation to the entropy condition.  Malformed descriptions
are rejected.  Independent repetition with an intermediate threshold
amplifies this constant gap to the standard parameters.
This proves the lemma without requiring a classical algorithm to find
the good support.  QCMA completeness requires existence of the classical
witness, not efficient discovery of it.

\paragraph{Attempt II: could spectral structure supply the hypothesis?}
The lemma suggests seeking a polynomial-description basis in which some
vector from the accepting spectral subspace has logarithmic measurement
entropy.  One might hope that the succinct circuit defining $A_x$ forces
this.  No such implication has been proved here.  Describing an operator
by a short circuit does not automatically describe one of its extremal
eigenvectors by a short state-preparation circuit.

Two examples prevent weaker surrogate statements from closing the gap.
For $A=|+\rangle\langle+|^{\otimes m}$, the top eigenvalue is one but
every computational-basis vector has acceptance $2^{-m}$.  A verifier
restricted to basis-state witnesses fails completely.  Nevertheless
$|+\rangle^{\otimes m}$ is prepared by $m$ Hadamard gates, so classical
verification is easy.  Large entropy in one chosen basis is therefore
not a witness-complexity lower bound.  In the Hadamard basis its entropy
is zero, exactly as the compression lemma allows.

Conversely, a counting argument shows that most states do not have
polynomial-size descriptions, but gives no efficiently verifiable
language for which \emph{all} QCMA protocols fail.  It must first link
those states to an explicit verifier's accepting subspace, and then
exclude alternative classical-witness verification strategies.  The
oracle separation source explicitly discusses this obstacle.
\sref{122}{1}

A spectral-gap projection does not solve the state-preparation issue
either.  Starting with a maximally mixed state and projecting onto an
accepting subspace of dimension $r$ among $2^m$ dimensions has success
probability $r/2^m$.  Even a perfectly implementable projection is not a
polynomial-time preparation when this probability is exponentially small.
Renormalizing after projection is postselection, not a free operation
in a QMA-to-QCMA simulation.

\paragraph{Stress test.}
The entropy hypothesis is only sufficient; it is not claimed to
characterize QCMA.  The polynomial exponent $C/\delta$ is fixed once
the language is fixed.  Letting it grow with $n$ would destroy the
complexity bound.  The preparation argument supplies an actual state
with a controlled trace-distance error and handles all classical
certificates on no-inputs, not only honest descriptions.

The existing oracle constructions are evidence that a fully relativizing
equality proof cannot succeed, but they are not unrelativized lower
bounds.  Nor does a separation between promise classes automatically
supply the language separation selected here.  Finally, an ordinary
QMA/QCMA separation would imply $\mathsf P\ne\mathsf{PSPACE}$, using
$\mathsf P\subseteq\mathsf{QCMA}\subseteq\mathsf{QMA}\subseteq
\mathsf{PSPACE}$; the exponential-dimensional acceptance operator can be
handled in polynomial space by the standard quantum-verification
simulation.  Thus making the last lower-bound step explicit would have
major additional consequences, not merely establish an encoding fact.

The companion script checks the rank-one acceptance example and the
trace-distance/truncation inequality on finite rational test vectors.
Those checks certify neither an entropy bound for all QMA languages
nor a complexity-class separation.

\paragraph{Outcome.}
\textbf{Outcome: Exploratory approach with unresolved gap.}

A precise accepting-witness entropy condition has been shown to imply
QCMA verification, with polynomial encoding size, preparation error,
and soundness handled explicitly.  No argument supplies that condition
for every QMA language; it is not necessary for every possible QCMA
protocol either.  The basis-state and postselection shortcuts fail for
specified reasons.  No equality, separation, or novelty claim is made.
% END RESEARCH ADDITION 122
\subsection{Sources}
\begin{itemize}\raggedright
\item[\textnormal{[S1]}]\hypertarget{source:122:1}{} John Bostanci, Jonas Haferkamp, Chinmay Nirkhe, and Mark Zhandry, Separating QMA from QCMA with a classical oracle, Introduction.
\url{https://eccc.weizmann.ac.il/report/2025/176/revision/1/download/}.
{\small\textit{Catalogue formulation source.}}
% BEGIN RESEARCH REFERENCES 122
\item[\textnormal{[E1]}]\hypertarget{extra:122:1}{} John Bostanci, Andrew Huang, and Vinod Vaikuntanathan.
\emph{Separating Quantum and Classical Advice with Good Codes}, ECCC TR26-020, revision 1, 12 April 2026. Abstract and introduction; the separation is relative to a classical oracle.
\url{https://eccc.weizmann.ac.il/report/2026/020/}.
{\small\textit{Consulted 27 September 2026 for the indicated result or scope.}}
% END RESEARCH REFERENCES 122
\end{itemize}


\end{document}
